\chapter{Systèmes similaires}
\label{chap:etat-art}

\section{Introduction}

Toute démarche d'ingénierie logicielle commence par une lecture attentive du terrain sur lequel la solution est appelée à s'inscrire. Le présent chapitre poursuit cette ambition~: nous examinons d'abord les enjeux de la mobilité intelligente et du transport inter-wilayas en Algérie, puis nous étudions individuellement les plateformes existantes de transport et de covoiturage afin de dégager les limites structurelles du marché et de positionner le projet \textit{RohWinBghit}.

\section{Contexte et enjeux de la mobilité inter-wilayas en Algérie}
\label{sec:presentation-domaine}

\subsection{Mobilité intelligente et connectivité mobile en Algérie}
\label{sec:mobilite-intelligente}

Le concept de \textit{Système de Transport Intelligent} (\textsc{its} --- \textit{Intelligent Transportation System}) désigne l'intégration de technologies de l'information et de la communication dans les infrastructures et les services de transport afin d'optimiser les flux de personnes et de marchandises. Dans le domaine spécifique de la mobilité partagée, cette intelligence se matérialise par la capacité de la plateforme à apparier dynamiquement l'offre et la demande en temps réel, à adapter la tarification aux conditions instantanées du marché, et à garantir la confiance entre inconnus par des mécanismes de vérification algorithmique.

L'efficacité d'un tel système repose de manière critique sur la qualité du réseau de télécommunications sous-jacent. En Algérie, l'infrastructure mobile a connu une évolution significative au cours des deux dernières décennies~: du \textsc{gprs} (2,5G) déployé au début des années~2000 vers le réseau \textsc{3g} inauguré par les opérateurs nationaux en 2013, puis la couverture \textsc{4g~lte} étendue progressivement à partir de 2016. Néanmoins, la réalité terrain impose une contrainte architecturale fondamentale~: sur les axes inter-wilayas traversant le Haut-Plateau et les régions steppiques, la connectivité mobile demeure intermittente. Les tunnels, les zones de faible densité de population et les transitions entre cellules de relais provoquent des coupures régulières. Cette réalité justifie scientifiquement l'adoption d'une architecture asynchrone et tolérante aux déconnexions~: requêtes idempotentes avec clé \texttt{Idempotency-Key}, files de messages persistantes (Bull/BullMQ), retentatives exponentielles via TanStack~Query, et persistance locale des données critiques sur le terminal mobile avant synchronisation différée.

\subsection{Transport collaboratif}
\label{sec:definition-covoiturage}

Le covoiturage (\textit{carpooling}) désigne le partage d'un véhicule privé entre un conducteur et un ou plusieurs passagers réalisant un trajet similaire, moyennant un partage des frais. Structuré dans les années~2000 avec l'émergence d'Internet puis du \textit{smartphone}~\cite{chan_carpooling_2012}, il a été industrialisé par des plateformes telles que \textit{BlaBlaCar} (France, 2006), qui ont facilité la mise en relation à grande échelle et introduit des mécanismes de confiance numérique.

Lorsqu'il est correctement encadré par une plateforme numérique, le covoiturage produit une triple valeur~: économique (mutualisation des frais, ~40\,\% d'économie par passager~\cite{blablacar_report_2023}), sociale (accessibilité aux zones mal desservies, démocratisation des déplacements pour les étudiants) et environnementale (réduction du nombre de véhicules, ~110~kg de CO\textsubscript{2} évités par passager/an~\cite{ademe_2022}).

\begin{figure}[!htbp]
\centering
\resizebox{\textwidth}{!}{%
\begin{tikzpicture}[
  node distance=2.2cm and 2.5cm,
  actor/.style={rectangle, rounded corners=4pt, draw=chapcolor, fill=rohlight, minimum width=3cm, minimum height=1cm, align=center, font=\footnotesize\bfseries\color{chapcolor}},
  platform/.style={rectangle, rounded corners=4pt, draw=rohaccent, fill=rohaccent!10, minimum width=3.5cm, minimum height=1.2cm, align=center, font=\footnotesize\bfseries\color{rohaccent}},
  arrow/.style={->, >=stealth, draw=chapcolor!80, thick, font=\scriptsize, align=center, text=black}
]
  % Nodes
  \node[actor] (passenger) {Passager\\(Demande / Paiement)};
  \node[platform, right=1.5cm of passenger] (platform) {\textbf{Plateforme RohWinBghit}\\\textit{Mise en relation, KYC, Séquestre}};
  \node[actor, right=1.5cm of platform] (driver) {Conducteur\\(Offre de places / Trajet)};

  % Arrows
  \path[arrow] (passenger) edge[bend left=15] node[above, pos=0.5] {1. Recherche \& Réservation} (platform);
  \path[arrow] (platform) edge[bend left=15] node[above, pos=0.4] {2. Notification de demande} (driver);
  \path[arrow] (driver) edge[bend right=38] node[above, pos=0.4] {3. Acceptation} (platform);
  \path[arrow] (platform) edge[bend left=15] node[below, pos=0.5] {4. Billet \& Code QR} (passenger);
  \path[arrow] (passenger) edge[bend right=42] node[below, pos=0.5, align=center] {5. Partage du trajet et Validation GPS} (driver);
  \path[arrow] (platform) edge[bend right=25] node[below, pos=0.4] {6. Versement fonds\\(Commission 12\,\% déduite)} (driver);
\end{tikzpicture}%
}
\caption{Schéma conceptuel de mise en relation et flux transactionnel dans RohWinBghit}
\label{fig:concept-covoiturage}
\end{figure}

\subsection{Enjeux spécifiques du contexte algérien}
\label{sec:transport-algerie}

L'Algérie présente une géographie singulière~: un territoire de plus de 2,38~millions de km\textsuperscript{2}, structuré en 69~wilayas, et marqué par une concentration urbaine sur le littoral. Les déplacements inter-wilayas restent fortement dépendants de la voiture individuelle et de l'autocar. Plusieurs travaux académiques nationaux se sont penchés sur la modélisation mathématique et l'optimisation des flux de covoiturage dynamique en Algérie, notamment à Béjaïa~\cite{bouzidi_covoiturage_dz_2022}, soulignant l'adéquation de ce mode de transport aux contraintes de mobilité locales. Quatre constats structurent le diagnostic.
\begin{enumerate}
  \item \textbf{Demande étudiante forte.} Les flux hebdomadaires entre les wilayas universitaires (Tlemcen, Oran, Constantine, Sétif, Alger) et les wilayas d'origine sont particulièrement intenses.
  \item \textbf{Offre saturée aux périodes de pointe.} Les autocars et taxis collectifs présentent des niveaux de confort et de ponctualité très inégaux.
  \item \textbf{Attractivité économique de la mutualisation.} Le coût total de possession d'un véhicule individuel est en hausse constante, rendant le partage des frais attractif.
  \item \textbf{Absence de plateforme nationale de référence.} Il n'existe pas de solution unifiée offrant une expérience de réservation longue distance de niveau international.
\end{enumerate}

\subsection{Marché numérique algérien}
\label{sec:opportunite-marche}

Le marché algérien de la mobilité numérique est émergent mais dynamique. Le taux de pénétration du \textit{smartphone} dépasse 75\,\%, et le pays comptait plus de 48~millions d'abonnements Internet mobile actifs en 2024~\cite{arpt_2024}. L'écosystème des paiements électroniques s'est structuré avec les cartes \textbf{CIB} et \textbf{Edahabia}, complétées par \textit{Baridi~Mob}~: plus de \textbf{7~millions d'opérations} de paiement électronique ont été enregistrées sur les premiers mois de 2025~\cite{aps_epaiement_2025}. Enfin, le pays compte plus de 1,7~million d'étudiants, dont une part importante se déplace régulièrement entre wilayas, constituant une base initiale naturelle pour une plateforme inter-wilayas.

\begin{figure}[!htbp]
\centering
\resizebox{\textwidth}{!}{%
\begin{tikzpicture}[xscale=1.5, yscale=0.07]
  % Grid
  \draw[very thin, gray!30, ystep=10, xstep=1] (0,0) grid (6,100);
  
  % Axes
  \draw[->, >=stealth, thick] (0,0) -- (6.5,0) node[right, font=\footnotesize] {Année};
  \draw[->, >=stealth, thick] (0,0) -- (0,105) node[above, font=\footnotesize] {Taux / Pénétration (\%)};

  % X-Ticks
  \draw (0,2) -- (0,-2) node[below, font=\tiny] {2020};
  \draw (1,2) -- (1,-2) node[below, font=\tiny] {2021};
  \draw (2,2) -- (2,-2) node[below, font=\tiny] {2022};
  \draw (3,2) -- (3,-2) node[below, font=\tiny] {2023};
  \draw (4,2) -- (4,-2) node[below, font=\tiny] {2024};
  \draw (5,2) -- (5,-2) node[below, font=\tiny] {2025};
  \draw (6,2) -- (6,-2) node[below, font=\tiny] {2026};

  % Y-Ticks
  \foreach \y in {0, 20, 40, 60, 80, 100}
    \draw (0.05,\y) -- (-0.05,\y) node[left, font=\tiny] {\y\%};

  % Data lines
  % Smartphone penetration
  \draw[ultra thick, chapcolor, mark=*] plot coordinates {
    (0, 52) (1, 58) (2, 63) (3, 69) (4, 75) (5, 81) (6, 85)
  };
  \node[chapcolor, right, font=\tiny\bfseries] at (6, 85) {Pénétration Smartphone};

  % Mobile Internet Active Users
  \draw[ultra thick, rohaccent, mark=square*] plot coordinates {
    (0, 38) (1, 41) (2, 43) (3, 46) (4, 48) (5, 52) (6, 55)
  };
  \node[rohaccent, right, font=\tiny\bfseries] at (6, 55) {Abonnements Internet Mobile (Millions)};

  % E-payment operations (scaled for representation or separate curve)
  \draw[ultra thick, green!60!black, mark=triangle*] plot coordinates {
    (0, 5) (1, 12) (2, 24) (3, 38) (4, 55) (5, 72) (6, 90)
  };
  \node[green!60!black, right, font=\tiny\bfseries] at (6, 90) {Transactions E-paiement (Index)};
\end{tikzpicture}%
}
\caption{Évolution des indicateurs du marché numérique algérien (2020--2026) --- Données combinées ARPT/APS}
\label{fig:market-trends}
\end{figure}

\section{BlaBlaCar}
\label{sec:blablacar}

BlaBlaCar (France, 2006) reste la référence européenne et internationale du covoiturage longue distance, avec une présence dans 22~pays et plus de 100~millions d'utilisateurs cumulés. Elle offre une structure complète de réservation, une messagerie interne et un système d'évaluation des membres. Son modèle de paiement (carte bancaire internationale obligatoire) et son absence d'implantation locale le rendent inopérant pour le marché national algérien, mais il constitue un modèle fonctionnel de référence.

\section{Yassir}
\label{sec:yassir}

Yassir est une application algérienne fondée en 2017, devenue le leader local du \gls{vtc} urbain et s'étendant à la livraison de repas et aux services financiers. Une étude empirique conduite à Ali~Menjeli (Constantine) a montré que son adoption a réduit le temps d'attente moyen et transformé les pratiques de mobilité~\cite{naim_yassir_2024}. 

D'un point de vue architectural, la plateforme Yassir ne propose aucune vérification biométrique de vivacité (\textit{liveness detection}). Son processus d'onboarding conducteur se limite à une vérification documentaire asynchrone, ce qui expose structurellement le système à des attaques par rejeu d'identité~: un acteur malveillant peut soumettre la photographie d'un tiers sans que le système ne détecte l'absence de présence physique. Par ailleurs, son modèle économique reste exclusivement urbain et de transport à la demande (VTC), ne couvrant aucun axe inter-wilayas de covoiturage régulier.

\paragraph{Heetch \& TemTem.} D'autres acteurs urbains algériens comme Heetch (introduit en 2020) et TemTem ont également consolidé le transport urbain à la demande. Ils partagent cependant les mêmes limites que Yassir : focalisation sur les trajets courts intra-wilayas et absence de modules de covoiturage longue distance ou de vérification biométrique KYC intégrée.

\section{Nroho}
\label{sec:nroho}

Nroho est une startup algérienne de covoiturage labellisée par le Ministère de l'Économie du Savoir et lauréate du 2\textsuperscript{e}~prix à l'\textit{IsDB Pitch Competition}. L'application revendique plus de 2~000~utilisateurs actifs et propose une vérification d'identité par scan de la \gls{cin} et \textit{selfie}, un \textit{matching} par affinités, un suivi GPS temps réel et un règlement en espèces. 

L'examen technique de son pipeline d'identité révèle cependant une vulnérabilité structurelle~: l'absence de réseau convolutif (\textsc{cnn}) dédié à la détection d'anti-\textit{spoofing} signifie qu'une photographie imprimée haute résolution présentée devant la caméra suffit à contourner le mécanisme de vérification. En outre, les transactions financières basées exclusivement sur un modèle fiduciaire physique induisent un risque de désintermédiation financière~: l'opérateur ne peut appliquer ni séquestre, ni frais de service garantîd sur les trajets ultérieurs négociés hors plateforme, impactant directement la viabilité du modèle économique.

\section{inDrive}
\label{sec:indrive}

inDrive est un acteur international classé application de mobilité partagée \#1 en Algérie au troisième trimestre 2025~\cite{sensortower_algeria_2025}. En décembre~2025, inDrive a étendu son offre vers le covoiturage inter-wilayas (service \textit{City~to~City}) sur une dizaine d'axes initiaux~\cite{indrive_interwilaya_2025}. 

La vérification d'identité y est purement documentaire (non biométrique), ce qui la rend vulnérable aux mêmes vecteurs d'attaque par photographie statique que les acteurs précédents. Plus fondamentalement, l'absence totale d'intégration des moyens de paiement électroniques algériens (\textsc{cib}, \textit{Edahabia}) force la totalité des transactions à s'effectuer en espèces, ce qui non seulement génère une opacité fiscale incompatible avec les exigences de la Loi~18-05 sur le commerce électronique~\cite{loi1805_2018}, mais empêche également la plateforme de mettre en œuvre un mécanisme de séquestre garantissant la protection financière du passager en cas d'annulation.

\section{Covoit'Dz (Université de Tlemcen)}
\label{sec:covoitdz}

Développée et soutenue en tant que projet académique au sein du Département d'Informatique de l'Université de Tlemcen~\cite{merad_covoitdz_2024}, l'application Android Covoit'Dz vise spécifiquement à faciliter la mise en relation pour le partage de trajets routiers en Algérie.
\begin{itemize}
  \item \textbf{Analyse de la contribution~:} Le projet démontre l'intérêt méthodologique de l'approche mobile Android pour structurer le marché national, en proposant une modélisation UML et une base de données relationnelle adaptée aux liaisons inter-wilayas.
  \item \textbf{Limitations scientifiques~:} Covoit'Dz souffre de l'absence totale de modules d'identité biométrique KYC. L'application repose sur un modèle déclaratif vulnérable, sans chiffrement des données personnelles au repos et sans algorithme de tarification dynamique. De plus, elle ne propose aucune intégration monétique avec l'écosystème interbancaire de la SATIM, limitant son champ d'action à un prototype informationnel.
\end{itemize}

\section{Synthèse des systèmes concurrents}
\label{sec:comparatif}

Afin de situer précisément notre contribution, nous présentons ci-dessous une matrice comparative des plateformes étudiées (Tableau~\ref{tab:comparatif}) suivie d'une analyse des limites du marché.

\begin{landscape}
\begin{table}[!htbp]
\centering
\footnotesize
\renewcommand{\arraystretch}{1.3}
\begin{tabularx}{\linewidth}{|l|c|c|c|c|c|X|}
\hline
\rowcolor{rohlight}
\textbf{Critère} & \textbf{BlaBlaCar} & \textbf{Yassir} & \textbf{Nroho} & \textbf{inDrive} & \textbf{Covoit'Dz} & \textbf{RohWinBghit} \\
\hline
Couverture en Algérie & Absent & Urbaine & Nationale & Nationale & Nationale & \textbf{Nationale (69 wilayas)} \\
\hline
Focus inter-wilayas & Oui (Europe) & Non & Oui & Oui (City to City) & Oui & \textbf{Oui (natif)} \\
\hline
Authentification KYC & Standard & Standard & CIN + \textit{selfie} & Documentaire & Néant & \textbf{Biométrique (Liveness + Anti-spoof)} \\
\hline
Seuils IA par rôle & Non & Non & Non & Non & Non & \textbf{Oui ($\geq 0{,}65$ / $\geq 0{,}60$)} \\
\hline
Paiement CIB / Edahabia & Non & Partiel & Non (cash) & Non (cash) & Non & \textbf{Sandbox + espèces} \\
\hline
Messagerie temps réel & Oui & Oui & Oui & Oui & Non & \textbf{Oui (Socket.IO)} \\
\hline
Tarification dynamique & Non & Oui & Non & Non & Fixe & \textbf{Oui ($m \in [1{,}0;\; 2{,}5]$)} \\
\hline
Chiffrement PII au repos & Non doc. & Non doc. & Non doc. & Non doc. & Non doc. & \textbf{AES-256-GCM} \\
\hline
Langues d'interface & FR / EN & AR / FR & AR / FR & AR / FR & FR & \textbf{AR / FR / EN} \\
\hline
\end{tabularx}
  \caption{Matrice comparative approfondie des plateformes de mobilité partagée --- État du marché au 1\textsuperscript{er} mai 2026}
  \label{tab:comparatif}
\end{table}

\begin{highlight}
\textbf{Synthèse du positionnement.} \textit{RohWinBghit} ambitionne d'être une plateforme algérienne de mobilité partagée inter-wilayas techniquement aboutie : vérification biométrique à trois niveaux (identité, vivacité, anti-\textit{spoofing}), méthodes de paiement locales (espèces, CIB, \textit{Edahabia}), couverture nationale et interface trilingue. Elle s'inscrit dans la continuité de l'écosystème existant (Yassir, Heetch, Nroho, inDrive) et y apporte une proposition spécifiquement pensée pour le segment longue distance.
\end{highlight}
\end{landscape}

\subsection{Limites du marché et positionnement de RohWinBghit}
\label{sec:limites}

L'analyse comparative met en évidence cinq limites structurelles majeures qui définissent autant d'axes de différenciation pour notre projet.

\begin{enumerate}[leftmargin=1.8em]
  \item \textbf{Confiance limitée (Lacune~1)~:} Les plateformes existantes se limitent à une vérification documentaire sans contrôle de vivacité ni anti-\textit{spoofing}. \textit{RohWinBghit} y répond par un pipeline biométrique à trois niveaux (OCR, comparaison faciale, détection de vivacité) avec des seuils différenciés par rôle.
  \item \textbf{Paiement électronique non intégré (Lacune~2)~:} Le règlement en espèces reste la norme quasi-exclusive. \textit{RohWinBghit} intègre trois méthodes coexistantes~: espèces (opérationnel), \textsc{cib} et \textit{Edahabia} (sandbox), avec une architecture extensible (patron \textit{Factory}).
  \item \textbf{Couverture interurbaine parcellaire (Lacune~3)~:} Les acteurs proposent des lignes géographiques rigides. \textit{RohWinBghit} prend en charge nativement les 69~wilayas avec interface trilingue arabe/français/anglais.
  \item \textbf{Absence d'outils de modération (Lacune~4)~:} Les applications manquent de panneaux de contrôle pour la validation \textsc{kyc} et la gestion des litiges. \textit{RohWinBghit} intègre une interface web d'administration dédiée.
  \item \textbf{Absence de tarification dynamique et de matching inversé (Lacune~5)~:} Il n'existe aucun moteur de matching intelligent ni mécanisme de \textit{surge pricing}. \textit{RohWinBghit} propose un prototype de tarification dynamique et un mécanisme de demande de trajet (matching inversé).
\end{enumerate}

Le projet ne prétend pas remplacer les acteurs existants, mais propose une réponse adaptée au contexte local sur le segment inter-wilayas. L'étude de Boureima~\textit{et~al.}~\cite{mounir_carpooling_developing_2022} identifie la confiance entre inconnus et la disponibilité du paiement électronique comme deux verrous majeurs de l'adoption du covoiturage en Afrique du Nord --- précisément les axes sur lesquels notre proposition se concentre. L'arrivée d'inDrive sur ce segment fin~2025~\cite{indrive_interwilaya_2025} et la présence durable de Nroho confirment l'existence d'une demande structurée. Le cadre réglementaire a été renforcé en 2025 par la \textbf{Loi~25-11}~\cite{loi2511_2025} sur la protection des données personnelles, qui concerne directement toute plateforme manipulant des données biométriques~\cite{anpdp_obligations_2026}.


\section{Spécification des exigences}
\label{sec:spec-exigences}

Cette section formalise la spécification des exigences fonctionnelles (tableau~\ref{tab:exigences-fonc}) ainsi que des exigences non fonctionnelles (détaillées sous forme de fiches techniques de l'ENF~1 à l'ENF~8 ci-dessous) qui régissent le comportement technique, sécuritaire et opérationnel de la plateforme.

\subsection{Exigences fonctionnelles}
\label{sec:recueil-besoins}
\label{sec:exigences-fonc}

La consolidation des besoins repose sur le croisement des trois sources d'analyse (questionnaire, étude comparative, entretiens). Les exigences les plus fréquemment citées concernent la confiance dans le conducteur, la visibilité du prix avant réservation, la disponibilité d'un canal de discussion direct et la prise en charge des moyens de paiement locaux.

\begin{table}[!htbp]
\centering
\small
\begin{tabularx}{\textwidth}{|l|l|X|l|l|}
\hline
\rowcolor{rohlight}
\textbf{ID} & \textbf{Priorité} & \textbf{Description courte} & \textbf{Acteur} & \textbf{UC} \\
\hline
EF1.1 & Haute & Inscription, authentification OTP et gestion profil & Passager, Conducteur & UC-AUTH \\
\hline
EF1.2 & Haute & Vérification d'identité biométrique (OCR + liveness) & Passager, Conducteur & UC-KYC \\
\hline
EF2.1 & Haute & Publication d'un trajet avec tarif dynamique & Conducteur & UC-PUB \\
\hline
EF2.2 & Moyenne & Validation ou refus des réservations & Conducteur & UC-BOOK \\
\hline
EF1.3 & Haute & Recherche de trajets avec filtres & Passager & UC-SEARCH \\
\hline
EF1.4 & Haute & Réservation et paiement (CIB/Edahabia/espèces) & Passager & UC-BOOK \\
\hline
EF1.5 & Haute & Messagerie temps réel et suivi GPS & Passager, Conducteur & UC-CHAT \\
\hline
EF3.1 & Haute & Tableau de bord, statistiques et gestion KYC manuel & Administrateur & UC-MOD \\
\hline
\end{tabularx}
  \caption{Exigences fonctionnelles principales}
  \label{tab:exigences-fonc}
\end{table}

\subsection{Exigences non fonctionnelles}
\label{sec:exigences-nonfonc}

Les exigences non fonctionnelles ont été dérivées de la revue systématique de Sogbe~\textit{et~al.}~\cite{sogbe_public_transport_2024} sur les déterminants de la satisfaction dans les systèmes de transport des pays en développement.

\begin{nonfuncreq}[title={\textbf{ENF~1 --- Sécurité}}]
Authentification forte (OTP + \gls{jwt}), chiffrement \gls{tls}~1.3 en transit, chiffrement au repos des documents \textsc{kyc} (AES-256-GCM), journalisation des actions administrateur.
\end{nonfuncreq}

\begin{nonfuncreq}[title={\textbf{ENF~2 --- Performance}}]
Temps de réponse API médian $\leq$ 200~ms, chargement initial $\leq$ 3~s sur réseau 4G.
\end{nonfuncreq}

\begin{nonfuncreq}[title={\textbf{ENF~3 --- Scalabilité}}]
Architecture \textit{back-end} sans état, cache Redis pour les requêtes fréquentes.
\end{nonfuncreq}

\begin{nonfuncreq}[title={\textbf{ENF~4 --- Disponibilité}}]
Objectif de \textit{uptime} mensuel de 99{,}5\,\%, sauvegardes quotidiennes.
\end{nonfuncreq}

\begin{nonfuncreq}[title={\textbf{ENF~5 --- Maintenabilité}}]
Couverture de tests unitaires $\geq$ 80\,\%, intégration continue, documentation Swagger.
\end{nonfuncreq}

\begin{nonfuncreq}[title={\textbf{ENF~6 --- Ergonomie et accessibilité}}]
Interface trilingue (arabe RTL, français, anglais) via \texttt{i18next}~; accessibilité \textsc{wcag}~2.1 niveau~A.
\end{nonfuncreq}

\begin{nonfuncreq}[title={\textbf{ENF~7 --- Conformité réglementaire}}]
Respect de la Loi~25-11~\cite{loi2511_2025} sur la protection des données personnelles.
\end{nonfuncreq}

\begin{nonfuncreq}[title={\textbf{ENF~8 --- Compatibilité multiplateforme}}]
Application native Android/iOS via React~Native et Expo.
\end{nonfuncreq}

\clearpage
\section{Modélisation des besoins (UML)}
\label{sec:analyse-fonctionnelle}

Cette section présente la modélisation fonctionnelle du système en identifiant ses acteurs (tableau~\ref{tab:acteurs}) et en modélisant leurs interactions sous la forme d'un diagramme global de cas d'utilisation (figure~\ref{fig:uc-global}). La cohérence de la spécification est renforcée par une matrice de traçabilité (tableau~\ref{tab:tracabilite}) et par la définition détaillée des récits utilisateurs (User Stories) présentés dans le tableau~\ref{tab:user-stories}.

\subsection{Identification des acteurs}
\label{sec:acteurs}

Le système met en relation trois acteurs humains et plusieurs acteurs techniques externes. Le tableau~\ref{tab:acteurs} récapitule leurs rôles respectifs.

\begin{table}[!htbp]
\centering
\begin{tabularx}{\textwidth}{|l|l|X|}
\hline
\rowcolor{rohlight}
\textbf{Acteur} & \textbf{Type d'acteur} & \textbf{Rôle principal} \\
\hline
\textbf{Passager} & Humain & Recherche un trajet, réserve, paie, communique, note le service. \\
\hline
\textbf{Conducteur} & Humain & Publie un trajet, gère les réservations, encaisse via le portefeuille. \\
\hline
\textbf{Administrateur} & Humain & Modère la plateforme, gère les litiges, effectue la revue manuelle du KYC. \\
\hline
\textbf{Service \textsc{kyc} (IA)} & Système (Interne) & Microservice Python FastAPI interne effectuant l'analyse biométrique. \\
\hline
\textbf{Passerelle paiement} & Système (Externe) & Traite les transactions \textsc{cib} et \textit{Edahabia} en environnement \textit{sandbox}. \\
\hline
\textbf{Service de notification}& Système (Externe) & Délivre les notifications \textit{push} (Firebase \textsc{fcm}) et \textsc{sms}. \\
\hline
\textbf{Service cartographie} & Système (Externe) & Fournit la géolocalisation et les itinéraires (Mapbox). \\
\hline
\end{tabularx}
  \caption{Acteurs du système et rôles associés}
  \label{tab:acteurs}
\end{table}

\subsection{Diagramme de cas d'utilisation global}
\label{sec:uc-global}

La figure~\ref{fig:uc-global} présente, sur une même vue, les trois acteurs humains et leurs cas d'utilisation principaux.

\begin{figure}[!htbp]
\centering
\includegraphics[width=0.95\linewidth,height=0.85\textheight,keepaspectratio]{fig_use_case.png}
\caption{Diagramme de cas d'utilisation consolidé (Passager, Conducteur, Administrateur)}
\label{fig:uc-global}
\end{figure}

Trois familles de cas se dégagent~: les cas \textit{transverses} (s'inscrire, se connecter par OTP, compléter le \textsc{kyc}), les cas \textit{transactionnels} (publier, rechercher, réserver, payer) et les cas \textit{relationnels} (messagerie, notation, signalement). Les relations \texttt{<<include>>} matérialisent les préconditions structurelles~; les relations \texttt{<<extend>>} illustrent les actions facultatives.

\subsection{Matrice de traçabilité exigences $\leftrightarrow$ cas d'utilisation}
\label{sec:matrice-tracabilite}

Le tableau~\ref{tab:tracabilite} établit la traçabilité entre les besoins fonctionnels et les cas d'utilisation.

\begin{table}[!htbp]
\centering
\footnotesize
\begin{tabularx}{\textwidth}{|X|l|l|}
\hline
\rowcolor{rohlight}
\textbf{Besoin fonctionnel} & \textbf{Acteur} & \textbf{Cas d'utilisation} \\
\hline
Inscription / authentification OTP & Passager, Conducteur & UC-AUTH \\
\hline
Vérification d'identité biométrique & Passager, Conducteur & UC-KYC \\
\hline
Publication d'un trajet & Conducteur & UC-PUB-TRIP \\
\hline
Recherche d'un trajet & Passager & UC-SEARCH \\
\hline
Réservation et paiement & Passager & UC-BOOKING \\
\hline
Communication temps réel & Passager, Conducteur & UC-CHAT \\
\hline
Notation et avis & Passager, Conducteur & UC-REVIEW \\
\hline
Modération et validation \textsc{kyc} & Administrateur & UC-MODERATION \\
\hline
\end{tabularx}
  \caption{Matrice de traçabilité besoins fonctionnels $\leftrightarrow$ cas d'utilisation}
  \label{tab:tracabilite}
\end{table}

\subsection{Spécification des User Stories}
\label{sec:user-stories}

Les exigences sont modélisées sous forme d'User Stories détaillées regroupées par rôle (tableau~\ref{tab:user-stories}). Chaque story respecte les critères \textbf{INVEST} (Indépendante, Négociable, Valeureuse, Estimable, Suffisamment petite, Testable).

Les récits utilisateurs associés aux différents rôles se déclinent ainsi :
\begin{itemize}
  \item \textbf{Conducteur} : Décrivent la publication de trajets inter-wilayas, la gestion des demandes de réservation et la configuration du profil du véhicule.
  \item \textbf{Passager} : Décrivent la recherche de trajets avec filtres avancés, la réservation avec paiement multi-méthodes et la réception de notifications.
  \item \textbf{Administrateur} : Décrivent la modération de la file d'attente KYC, la validation manuelle et la gestion générale des comptes utilisateurs.
\end{itemize}

\begin{table}[!htbp]
\centering
\small
\begin{tabularx}{\textwidth}{|l|l|X|X|l|}
\hline
\rowcolor{rohlight}
\textbf{ID} & \textbf{Acteur} & \textbf{Description (En tant que... je veux... afin de...)} & \textbf{Critères d'acceptation} & \textbf{Priorité} \\
\hline
US-01 & Conducteur & En tant que conducteur vérifié, je veux publier un trajet inter-wilayas avec les wilayas de départ/arrivée, l'heure et le prix, afin d'attirer des passagers et diviser mes frais. & - Sélection obligatoire des 69 wilayas.\newline - Véhicule pré-enregistré valide.\newline - Prix dans les limites recommandées. & Haute \\
\hline
US-02 & Conducteur & En tant que conducteur, je veux pouvoir accepter ou rejeter chaque demande de réservation, afin de contrôler la composition de mon équipage. & - Notification instantanée.\newline - Actions d'acceptation/refus.\newline - Libération automatique des places en cas de refus. & Haute \\
\hline
US-03 & Passager & En tant que passager, je veux rechercher un trajet par wilaya et date, avec des filtres avancés, afin de planifier mon voyage. & - Formulaire trilingue.\newline - Tri par prix, note, heure.\newline - Statut KYC du conducteur visible. & Haute \\
\hline
US-04 & Passager & En tant que passager vérifié, je veux réserver une place et payer via Edahabia, CIB ou espèces, afin de sécuriser mon trajet. & - Intégration sandbox SATIM.\newline - Notification push conducteur.\newline - Remboursement auto si refus. & Haute \\
\hline
US-05 & Admin & En tant qu'administrateur, je veux consulter et traiter la file d'attente KYC, afin de garantir la sécurité. & - Comparaison visuelle des documents.\newline - Scores InsightFace/Liveness.\newline - Notification auto de décision. & Critique \\
\hline
\end{tabularx}
\caption{User Stories détaillées et critères d'acceptation}
\label{tab:user-stories}
\end{table}

\clearpage
\subsection{Fiches descriptives et diagrammes de séquences des cas d'utilisation clés}
\label{sec:cu-textuel}

Cette section fournit les fiches descriptives détaillées des cas d'utilisation les plus critiques du système, accompagnées de leurs diagrammes de séquences systèmes correspondants (DSS) pour illustrer les interactions dynamiques entre les acteurs et la plateforme.

\subsubsection{UC-AUTH --- Inscription et authentification}
\label{sec:uc-auth}

\begin{itemize}
  \item \textbf{Acteur principal :} Utilisateur non connecté.
  \item \textbf{Pré-conditions :} Numéro de téléphone algérien valide.
  \item \textbf{Scénario nominal :}
    \begin{enumerate}
      \item L'utilisateur saisit son numéro de téléphone.
      \item Le système envoie un code OTP par SMS.
      \item L'utilisateur valide le code OTP.
      \item L'utilisateur complète les détails de son profil.
      \item Le système crée le compte et donne accès à la plateforme.
    \end{enumerate}
  \item \textbf{Post-conditions :} Compte créé, jeton \gls{jwt} émis, statut \textsc{kyc} initialisé à \textit{unverified}.
\end{itemize}

L'inscription repose sur une vérification du numéro de téléphone par code OTP, suivie de la saisie des informations personnelles (figure~\ref{fig:seq-inscription}). Le processus de connexion ultérieur par code OTP est modélisé dans la figure~\ref{fig:seq-auth}.

\begin{figure}[H]
\centering
\includegraphics[width=0.85\linewidth]{fig_sequence_inscription.png}
\caption{Séquence --- Inscription d'un nouvel utilisateur}
\label{fig:seq-inscription}
\end{figure}

\begin{figure}[H]
\centering
\includegraphics[width=0.85\linewidth]{fig_sequence_authentification.png}
\caption{Séquence --- Authentification et connexion (OTP)}
\label{fig:seq-auth}
\end{figure}

\subsubsection{UC-KYC --- Vérification d'identité biométrique}
\label{sec:uc-kyc}

\begin{itemize}
  \item \textbf{Acteur principal :} Utilisateur authentifié.
  \item \textbf{Pré-conditions :} Compte utilisateur actif.
  \item \textbf{Scénario nominal :}
    \begin{enumerate}
      \item L'utilisateur téléverse le recto et le verso de sa carte nationale d'identité (CIN).
      \item L'utilisateur effectue une capture vidéo de vivacité (3 secondes).
      \item Le microservice d'intelligence artificielle procède à l'évaluation asynchrone des descripteurs biométriques à l'aide de réseaux de neurones convolutifs (ArcFace) et d'un pipeline de détection de vivacité.
      \item Le système rend une décision automatique (verified/pending\_review/rejected).
    \end{enumerate}
  \item \textbf{Post-conditions :} Statut \textsc{kyc} mis à jour dans le profil utilisateur.
\end{itemize}

Le flux \textsc{kyc} (figure~\ref{fig:seq-kyc}) constitue le scénario le plus critique~: téléversement du document, capture vidéo de vivacité, traitement asynchrone par le pipeline IA (comparaison faciale InsightFace, détection de vivacité, anti-\textit{spoofing}), décision automatique à trois états (\textit{verified}, \textit{pending\_review}, \textit{rejected}), et éventuelle revue manuelle.

\begin{figure}[H]
\centering
\includegraphics[width=0.95\linewidth]{fig_sequence_kyc.png}
\caption{Séquence --- Vérification \textsc{kyc} biométrique (3~états de décision)}
\label{fig:seq-kyc}
\end{figure}

\subsubsection{UC-PUB-TRIP --- Publication et matching de trajets}
\label{sec:uc-pub-trip}

\begin{itemize}
  \item \textbf{Acteur principal :} Conducteur vérifié.
  \item \textbf{Pré-conditions :} Statut KYC \textit{verified}, véhicule enregistré.
  \item \textbf{Scénario nominal :}
    \begin{enumerate}
      \item Le conducteur sélectionne les wilayas de départ et d'arrivée.
      \item Il choisit son véhicule pré-enregistré.
      \item Il saisit la date, l'heure, le prix par place et le nombre de places disponibles.
      \item Le système valide et publie le trajet.
    \end{enumerate}
  \item \textbf{Post-conditions :} Trajet visible dans les résultats de recherche.
\end{itemize}

\begin{figure}[H]
\centering
\includegraphics[width=0.95\linewidth]{fig_sequence_trip_request.png}
\caption{Séquence --- Demande de trajet (\textit{matching} passager $\to$ conducteurs)}
\label{fig:seq-trip-request}
\end{figure}

\subsubsection{UC-BOOKING --- Réservation et paiement}
\label{sec:uc-booking}

\begin{itemize}
  \item \textbf{Acteur principal :} Passager vérifié.
  \item \textbf{Pré-conditions :} Statut KYC \textit{verified}, trajet disponible sélectionné.
  \item \textbf{Scénario nominal :}
    \begin{enumerate}
      \item Le passager sélectionne le trajet et le nombre de places.
      \item Il choisit le mode de paiement (carte CIB/Edahabia ou espèces).
      \item Le système vérifie la disponibilité des places et calcule le prix (FareService).
      \item Le passager valide le paiement via la passerelle (SATIM/Edahabia sandbox).
      \item Le système confirme la réservation et notifie le conducteur.
    \end{enumerate}
  \item \textbf{Post-conditions :} Réservation créée à l'état \textit{confirmed}, places débitées, notification envoyée.
\end{itemize}

\begin{figure}[H]
\centering
\includegraphics[width=0.95\linewidth]{fig_sequence_booking.png}
\caption{Séquence --- Réservation avec paiement multi-méthodes (vue technique)}
\label{fig:seq-reservation}
\end{figure}

\subsubsection{UC-CHAT et UC-MODERATION --- Messagerie et modération administrative}
\label{sec:uc-chat-mod}

\begin{itemize}
  \item \textbf{UC-CHAT --- Messagerie interne en temps réel :}
    \begin{itemize}
      \item \textit{Acteurs :} Passager et Conducteur.
      \item \textit{Pré-conditions :} Réservation confirmée active.
      \item \textit{Scénario nominal :} Messages de type texte via Socket.IO.
     \end{itemize}
  \item \textbf{UC-MODERATION --- Modération et revue KYC manuelle :}
    \begin{itemize}
      \item \textit{Acteur :} Administrateur.
      \item \textit{Scénario nominal :} Consultation de la file \textit{pending\_review}, examen visuel comparatif des documents, validation ou rejet avec motif, mise à jour du statut.
    \end{itemize}
\end{itemize}

\begin{figure}[H]
\centering
\includegraphics[width=0.8\linewidth]{fig_kyc_review.png}
\caption{Flux d'activité de la modération et revue manuelle du KYC par l'administrateur}
\label{fig:seq-kyc-review}
\end{figure}

\subsection{Modélisation du comportement global (Diagramme d'activité)}
\label{sec:activite-global}

Le parcours global d'un utilisateur, depuis l'authentification jusqu'à la fin du trajet, est représenté par le diagramme d'activité (figure~\ref{fig:activite-booking}).

\begin{figure}[!htbp]
\centering
\includegraphics[width=0.85\linewidth,height=0.85\textheight,keepaspectratio]{fig_passenger_flow.png}
\caption{Diagramme d'activité --- parcours global de la plateforme}
\label{fig:activite-booking}
\end{figure}

Après authentification et \textsc{kyc}, l'utilisateur est orienté vers l'une des deux branches selon son rôle. Les n\oeuds de décision matérialisent les ruptures possibles (places insuffisantes, paiement échoué), chacune associée à une compensation (libération de places, remboursement).

\section{Conclusion}

Ce chapitre a replacé le projet \textit{RohWinBghit} dans son contexte algérien en dressant l'état de l'art du transport partagé, en analysant les solutions existantes et en formalisant l'ensemble des exigences fonctionnelles, non fonctionnelles et la modélisation UML des besoins du système. L'étude comparative met en évidence cinq lacunes structurelles du marché et justifie les choix de conception décrits au Chapitre~II.
